\scanpage{120}
\uline{Proposition} $X$ préschéma loc. noethérien \figure{\begin{tikzpicture}\node[draw,circle,inner sep=2pt]{\shortstack{?\\\uline{réduit}}};\end{tikzpicture}}.
\margin{[il suffit $X$ réduit aux pts de $Y$] Car la \unsure{condition} est locale}

$f\in\Gamma(X,\underline{\mathcal O}_X)$, $f$ non diviseur de 0; $Y=\deleted{\unsure{X_f}}V(f)$.

(i) $\underline{\mathcal O}_X/f\underline{\mathcal O}_X$ réduit

\deleted{(ii)} \inserted{($\alpha$) On suppose} les $\underline{\mathcal O}_x$ sont des ``bons'' anneaux, i.e. leurs \unsure{quotients} intègres satisfont i) Condition de finitude ii) Condition des chaînes stricte

Soient $Y_i$ les composantes irréductibles de $Y$.

Soit $Z$ \uline{affine} sur $X$, et $g\deleted{\unsure{\xi}}$ une section rationnelle de $Z$ sur $X$, telle que

(i) $g$ définie dans $X_f$ (ii) $g$ définie en les pts génériques $y_i$ des $Y_i$.

\uline{Alors $g$ est définie partout}
\margin{contre-exemple si $Z$ n'est pas affine ?? $X$ \uline{réduit} est-il essentiel ??}

On est \unsure{ramené au cas où} $X=\operatorname{Spec}(A)$, $A$ local, et où $Z=X[T]$, i.e. $g\in K$ corps des fractions de $A$.

\uline{Lemme 1} Soit $Y$ une partie fermée de $X$ loc. noeth. réduit satisfaisant ($\alpha$), \deleted{$Z$ affine sur $X$} $g$ une \deleted{section} fonction rationnelle de $X$, définie dans $X-Y$ et les points génériques des
\margin{119}
